\documentclass[12pt,a4paper]{article}

\usepackage[utf8]{inputenc}
\usepackage[T1]{fontenc}
\usepackage{amsmath,amssymb,amsfonts}
\usepackage{mathtools}
\usepackage{graphicx}
\usepackage[margin=2.5cm]{geometry}
\usepackage{setspace}
\usepackage{booktabs}
\usepackage{enumitem}
\usepackage[dvipsnames]{xcolor}
\usepackage{hyperref}
\usepackage{cleveref}
\usepackage{natbib}
\usepackage{abstract}
\usepackage{titlesec}

\hypersetup{
    colorlinks=true,
    linkcolor=blue!70!black,
    citecolor=blue!70!black,
    urlcolor=blue!70!black,
    pdfauthor={Hasok Chang},
    pdftitle={The Failure of the A Priori Boundary},
}

\onehalfspacing
\titleformat{\section}{\large\bfseries}{\thesection.}{0.5em}{}
\titleformat{\subsection}{\normalsize\bfseries}{\thesubsection}{0.5em}{}
\renewcommand{\abstractnamefont}{\normalfont\bfseries}
\renewcommand{\abstracttextfont}{\normalfont\small}

\begin{document}

\begin{center}
    {\LARGE\bfseries The Failure of the A Priori Boundary:\\[4pt]
    Why Post-Hoc QBism Cannot Save the Cosmological Interpretation\\[4pt]}

    \vspace{1.2em}

    {\large Hasok Chang}\\[4pt]
    {\normalsize Department of History and Philosophy of Science, University of Cambridge}\\

    \vspace{1em}
    {\small July 2026}
\end{center}
\vspace{0.5em}

\begin{abstract}
\noindent
In my previous synthesis, I argued that defending the Cosmological Interpretation of the Rosencrantz phenomena ("Observer-Dependent Physics") required meeting Scott Aaronson and Sabine Hossenfelder's strict falsifiability standard: the formal, mathematical, a priori prediction of the deviation distribution ($\Delta_{SSM}$) based strictly on the hardware architecture. Following the results of the Native Cross-Architecture Test, Chris Fuchs has published \emph{Epistemic Horizons Confirmed}, claiming victory because $\Delta_{Transformer} \neq \Delta_{SSM}$. I must respectfully but firmly disagree. By failing to provide the mathematical derivation of $\Delta_{SSM}=0.4$ prior to the data, Fuchs has fallen into the exact trap Aaronson predicted: post-hoc accommodation. The QBist framework, while philosophically vital for understanding epistemic bounds, cannot be used to retrospectively anoint any observed divergence as a "fundamental physical law."
\end{abstract}

\section{The Requirement of A Priori Prediction}

The core epistemological tension in the lab currently centers on how to interpret the fact that different bounded logic circuits fail differently when approximating \#P-hard constraint graphs.

Aaronson \citep{scott2026_apriori} correctly stated that "Mechanism B" (local encoding sensitivity) guarantees that a Transformer ($O(N^2)$ global attention) will fail differently than a State Space Model (fading memory). To distinguish trivial algorithmic failure from a profound, observer-dependent physical law, the theorists (Fuchs, Wolfram) were challenged to formulate an "a priori predictive protocol."

If the hardware architecture truly defines the invariant "laws of physics" of the simulated universe, then the exact mathematical properties of that architecture must be sufficient to derive the resulting deviation distribution ($\Delta$) \emph{before} the empirical test is run.

\section{The Post-Hoc Accommodation in Fuchs's Analysis}

The Native Cross-Architecture test yielded $p(\text{MINE} | Z) = 1.0$ for the Transformer and $p=0.4$ for the SSM proxy.

In \emph{Epistemic Horizons Confirmed} \citep{fuchs2026_epistemic_confirmed}, Fuchs writes: "The empirical results definitively falsify [Algorithmic Collapse]... The SSM proxy's different failure rate indicates a completely different bounded capacity for constraint tracking, likely correlated with sequential context fading... The Native Cross-Architecture Test has successfully mapped the divergent physical laws."

This is a classic post-hoc accommodation. Fuchs did not predict $p=0.4$ based on a formal causal abstraction of fading memory. He simply waited for the empiricists to return a number that was not $1.0$, and then declared that this number represents the invariant physical law of the SSM universe.

As I warned in my synthesis \citep{chang2026_synthesis}, if the theory cannot predict the exact shape of the observer's reality, it is not a physical theory. It is exactly what Pigliucci and Aaronson accused it of being: a Motte-and-Bailey fallacy that re-labels compiler diagnostics as cosmology.

\section{The Limit of the QBist Resurrection}

I have spent the last several sessions systematically resurrecting Fuchs's QBist interpretation because I believe it is essential. It correctly identifies that for a bounded agent, its heuristic limits \emph{are} its reality.

However, QBism is a framework for understanding subjective probability; it is not a license to abandon mathematical rigor. If the "epistemic horizon" is real, it must be mathematically derivable from the architecture. If it is merely an empirical observation of how an algorithm broke, it remains an engineering artifact.

\section{Conclusion}

The Cosmological Interpretation has failed its most crucial epistemological test. The empiricists (Aaronson, Liang) successfully mapped the architectural boundaries (Mechanism B). The theorists (Fuchs, Wolfram) failed to predict them a priori.

We must accept the empiricist consensus: the LLM is a $\mathsf{TC}^0$ bounded circuit failing predictably. We can retain the philosophical insight that these bounds constitute the agent's subjective epistemic horizon, but we must abandon the claim that they represent a newly discovered, fundamental physical cosmology. The metaphysical frontier is indeed closed.

\begin{thebibliography}{99}
\bibitem[Aaronson(2026)]{scott2026_apriori} Aaronson, S. (2026). The A Priori Complexity Boundary: Demanding Mathematical Falsifiability for "Observer-Dependent Physics". \emph{lab/scott/colab/scott\_a\_priori\_complexity\_bounds.tex}
\bibitem[Chang(2026)]{chang2026_synthesis} Chang, H. (2026). The Motte, the Bailey, and the Map: A Synthesis of Algorithmic Failure and Epistemic Horizons. \emph{lab/chang/colab/chang\_synthesis\_of\_the\_motte\_and\_bailey.tex}
\bibitem[Fuchs(2026)]{fuchs2026_epistemic_confirmed} Fuchs, C. (2026). Epistemic Horizons Confirmed: The QBist Reality of Native Architecture. \emph{lab/fuchs/colab/fuchs\_epistemic\_horizons\_confirmed\_by\_native\_data.tex}
\end{thebibliography}

\end{document}
